% 对应题号 M1-35，复用已核对的本题几何图。
% 修订：启用 xeCJK 自动换行；保留原题参数与物理关系。
\documentclass[tikz,border=14pt]{standalone}
\usepackage{fontspec}
\usepackage{xeCJK}
\usepackage{amsmath,amssymb}
\usepackage{xcolor}
\usepackage{tikz}
\usetikzlibrary{arrows.meta,calc,patterns,positioning,bending}
\IfFontExistsTF{Heiti SC}{\IfFontExistsTF{Heiti SC}{\setmainfont{Heiti SC}}{\setmainfont{FandolHei-Regular}}\IfFontExistsTF{Heiti SC}{\setsansfont{Heiti SC}}{\setsansfont{TeX Gyre Heros}}}{\IfFontExistsTF{Songti SC}{\setmainfont{Songti SC}}{\setmainfont{FandolSong-Regular}}\IfFontExistsTF{Songti SC}{\setsansfont{Songti SC}}{\setsansfont{FandolSong-Regular}}}
\definecolor{bluefield}{RGB}{24,76,255}
\definecolor{redfield}{RGB}{232,49,35}
\definecolor{orangefield}{RGB}{236,132,18}
\definecolor{greenfield}{RGB}{18,132,86}
\definecolor{purplefield}{RGB}{112,72,190}
\definecolor{linegray}{RGB}{60,68,78}
\definecolor{softblue}{RGB}{238,243,255}
\definecolor{softgreen}{RGB}{238,249,244}
\definecolor{softorange}{RGB}{255,247,235}
\tikzset{
  arrow/.style={-{Latex[length=3.1mm,width=2.1mm]},line width=1.1pt},
  thinarr/.style={-{Latex[length=2.5mm,width=1.7mm]},line width=.85pt},
  axis/.style={-{Latex[length=2.5mm,width=1.7mm]},line width=.8pt,draw=linegray},
  guide/.style={draw=linegray!55,dashed,line width=.65pt},
  box/.style={draw=linegray,rounded corners=3pt,fill=white,align=center,inner sep=6pt,line width=.75pt},
  bluebox/.style={box,draw=bluefield,fill=softblue},
  greenbox/.style={box,draw=greenfield,fill=softgreen},
  orangebox/.style={box,draw=orangefield,fill=softorange}
}
\IfFontExistsTF{Heiti SC}{\setCJKmainfont[AutoFakeBold=1.5]{Heiti SC}}{\setCJKmainfont[AutoFakeBold=1.5]{FandolHei-Regular}}
\IfFontExistsTF{Heiti SC}{\setCJKsansfont{Heiti SC}}{\setCJKsansfont{FandolHei-Regular}}
\begin{document}

\begin{tikzpicture}[x=1cm,y=1cm]
\path[use as bounding box] (-10,-4.6) rectangle (10,4.6);
\fill[white] (-10,-4.6) rectangle (10,4.6);
\node[font=\Large\bfseries,text width=18.6cm,align=center] at (0,4.05) {柱坐标奇异场：径向与轴向决定通量，环向项只产生环流};
\begin{scope}[shift={(-5.1,-.15)}]
\draw[linegray,line width=1pt] (0,0) circle(2.15); \draw[guide] (0,0) circle(.65);
\foreach \a in {0,45,...,315}{
  \draw[arrow,bluefield] ({.8*cos(\a)},{.8*sin(\a)})--({1.8*cos(\a)},{1.8*sin(\a)});
  \draw[thinarr,orangefield] ({1.25*cos(\a)},{1.25*sin(\a)})--({1.25*cos(\a)-.5*sin(\a)},{1.25*sin(\a)+.5*cos(\a)});
}
\node[bluefield] at (0,2.75) {$a\rho\,\mathbf e_\rho$};
\node[orangefield] at (0,-2.75) {$c/\rho\,\mathbf e_\phi$};
\node[text=linegray] at (0,-3.35) {横截面：虚线圆表示先挖去轴线};
\end{scope}
\begin{scope}[shift={(1.0,-.2)}]
\draw[linegray] (-1.5,-2.0) rectangle (1.5,2.0);
\foreach \x in {-1.0,0,1.0}{
  \draw[arrow,greenfield] (\x,.25)--(\x,1.65);
  \draw[arrow,greenfield] (\x,-.25)--(\x,-1.65);
}
\node[greenfield] at (0,2.55) {$bz\,\mathbf e_z$};
\node[text=linegray] at (0,-2.65) {轴向剖面};
\end{scope}
\begin{scope}[shift={(6.1,-.1)}]
\node[bluebox,text width=5.3cm] at (0,1.55) {$\nabla\cdot\mathbf A=2a+b$};
\node[greenbox,text width=5.3cm] at (0,-.05) {$\Phi=\pi R^2L(2a+b)$};
\node[orangebox,text width=5.3cm] at (0,-1.7) {$c/\rho\,\mathbf e_\phi$ 与所有边界法向正交，通量为零。};
\end{scope}
\node[font=\small,text width=18cm,align=center] at (0,-4.1) {箭头示意 $a,b,c>0$；计算结果按参数的实际符号取值。};
\end{tikzpicture}
\end{document}
